% Part: first-order-logic
% Chapter: completeness
% Section: downward-ls

\documentclass[../../../include/open-logic-section]{subfiles}

\begin{document}

\olfileid{fol}{com}{dls}
\olsection{قضیهٔ لوونهایم--اسکولم}

قضیهٔ لوونهایم--اسکولم می‌گوید اگر نظریه‌ای مدلی نامتناهی داشته باشد،
مدلی نیز دارد که حداکثر !!{denumerable} است. یک پیامد بی‌درنگ این واقعیت
آن است که منطق مرتبهٔ اول نمی‌تواند بیان کند که اندازهٔ !!a{structure}
!!{nonenumerable} است: هر !!{sentence} یا مجموعه‌ای از !!{sentence}s که
در همهٔ !!{structure}s !!{nonenumerable} ارضا شود، در یک ساختار
!!{enumerable} نیز ارضا می‌شود.

\begin{thm}
\ollabel{thm:downward-ls} اگر $\Gamma$ سازگار باشد، !!a{enumerable}
مدل دارد؛ یعنی در !!a{structure} ارضا می‌شود که دامنه‌اش یا متناهی است
یا !!{denumerable}.
\end{thm}

\begin{proof}
اگر $\Gamma$ سازگار باشد، !!{structure}~$\Struct M$ که از اثبات قضیهٔ
تمامیت به دست می‌آید دامنه‌ای چون $\Domain{M}$ دارد که اندازه‌اش از
مجموعهٔ ترم‌های بستهٔ زبان هنکینی~$\Lang L'$ بزرگ‌تر نیست. پس
$\Struct M$ حداکثر !!{denumerable} است.
\end{proof}

\begin{thm}
\ollabel{noidentity-ls} اگر $\Gamma$ مجموعه‌ای سازگار از !!{sentence}s
در زبان منطق مرتبهٔ اول بدون همانی باشد، مدلی !!a{denumerable} دارد؛
یعنی در !!a{structure} ارضا می‌شود که دامنه‌اش نامتناهی و
!!{enumerable} است.
\end{thm}

\begin{proof}
اگر $\Gamma$ سازگار باشد و هیچ جمله‌ای شامل همانی نداشته باشد،
!!{structure}~$\Struct M$ که از اثبات قضیهٔ تمامیت به دست می‌آید دامنه‌ای
چون $\Domain{M}$ دارد که عیناً مجموعهٔ ترم‌های بستهٔ زبان~$\Lang L'$
است. پس $\Struct{M}$ !!{denumerable} است، زیرا $\Trm[L']$ چنین است.
\end{proof}

\begin{ex}[پارادوکس اسکولم]
نظریهٔ مجموعه‌های زرملو--فرانکل~$\Log{ZFC}$ چارچوبی بسیار نیرومند است
که تقریباً همهٔ گزاره‌های ریاضی، از جمله واقعیت‌هایی دربارهٔ اندازهٔ
مجموعه‌ها، را می‌توان در آن بیان کرد. برای نمونه، $\Log{ZFC}$ می‌تواند
اثبات کند که مجموعهٔ~$\Real$ اعداد حقیقی !!{nonenumerable} است؛ می‌تواند
قضیهٔ کانتور را اثبات کند که مجموعهٔ توانی هر مجموعه از خود آن مجموعه
بزرگ‌تر است؛ و جز این‌ها. اگر $\Log{ZFC}$ سازگار باشد، همهٔ مدل‌هایش
نامتناهی‌اند و افزون بر این، همگی شامل !!{element}sی هستند که نظریه
دربارهٔ آن‌ها می‌گوید !!{nonenumerable}اند؛ مانند عنصری که قضیهٔ
$\Log{ZFC}$ دربارهٔ وجود مجموعهٔ توانی اعداد طبیعی را صادق می‌سازد.
بنا بر قضیهٔ لوونهایم--اسکولم، $\Log{ZFC}$ مدل‌های !!{enumerable} نیز
دارد---مدل‌هایی که شامل مجموعه‌های «!!{nonenumerable}»اند اما خود
!!{enumerable} هستند.
\end{ex}

\end{document}
